% Source-derived chapter 04: Hodge-theoretic preliminaries
% Original source: canonical-representations-surface-groups
\chapter{Hodge-theoretic preliminaries}
\label{canonical-representations-surface-groups-chapter-04}
\source{canonical-representations-surface-groups}

\begin{remark}[Source section]
\label{canonical-representations-surface-groups-chapter-04-source}
\source{canonical-representations-surface-groups}
This chapter records the source section; no Lean declaration has yet been checked for this material.
\notready
\end{remark}

% The source section title was: Hodge-theoretic preliminaries
\label{section:Hodge-theoretic-preliminaries}
\source{canonical-representations-surface-groups}
In this section we recall some preliminaries on polarizable complex variations
of Hodge structure (complex PVHS), mixed Hodge theory, and Simpson-Mochizuki's non-abelian Hodge theory.

\subsection{Variations of mixed Hodge structure}
\label{subsection:vmhs}
\source{canonical-representations-surface-groups}
Our main goal in this section is \autoref{theorem:unitary-MHS},
which describes the variation of mixed Hodge structure on the cohomology of a
unitary local system on a family of punctured curves.
We then describe the resulting bigrading explicitly in
\autoref{lemma:three-part-hodge-structure}.


We now review pertinent notation.
We refer the reader to \cite[\S5]{LL:geometric-local-systems} for basic definitions
related to complex PVHS.
Let $\mathbb V$ be a unitary local system on a smooth quasi-projective variety.
We say that \emph{$\mathbb{V}$ has a real structure} if there exists a real orthogonal local system $\mathbb{V}_\mathbb{R}$ such that
$\mathbb{V}\simeq \mathbb{V}_\mathbb{R}\otimes_\mathbb{R}\mathbb{C}$. Note that
for any unitary local system $\mathbb{V}$, the local system
$\mathbb{V}\oplus\mathbb{V}^\vee$ has a natural real structure.

Let $\pi: \mathscr{C} \to \mathscr{M}$ be a family of $n$-pointed curves as in
\autoref{notation:versal-family}.
Further assume that $\mathscr{M}$ is smooth.
Let $\pi^\circ: \mathscr{C}^\circ\to
\mathscr{M}$ be the associated family of punctured curves,
$j : \mathscr C^\circ \to \mathscr C$ the inclusion,
$\mathscr{D}=\mathscr{C}\setminus\mathscr{C}^\circ$,
and $\mathbb V$ a unitary local system on $\mathscr C^\circ$.
Below we will use the notion of the Deligne canonical extension of a flat vector bundle from $\mathscr C^\circ$
to $\mathscr C$, as described in
\cite[Remarques 5.5(i)]{deligne:regular-singular} or \cite[Definition
4.1.2]{LL:geometric-local-systems}.

See \cite[Definitions 14.44, 14.45, and 14.49]{peters2008mixed} for the definitions of variations of mixed Hodge structure, graded-polarizability
and
admissibility in the rational case; the real case is analogous and discussed in \cite[Definition 3.4]{steenbrink1985variation}.

The following is well-known but perhaps difficult to extract from the
literature.
\begin{theorem}
\source{canonical-representations-surface-groups}
\notready\label{theorem:unitary-MHS}
	Suppose $\mathbb{V}$ is a unitary local system on $\mathscr{C}^\circ$
	with real structure $\mathbb{V}_\mathbb{R}$. Then,
	$R^1\pi^\circ_*\mathbb{V}_\mathbb{R}$ underlies an admissible
	graded-polarizable real variation of mixed Hodge structure with weights
	in $[1,2]$, and weight filtration given by
$$W^1R^1\pi^\circ_*\mathbb{V}_\mathbb{R}:=R^1\pi_*j_*\mathbb{V}_\mathbb{R}\subset R^1\pi^\circ_*\mathbb{V}_\mathbb{R}=:W^2R^1\pi^\circ_*\mathbb{V}_\mathbb{R}.$$
Let $(\mathscr{E}, {\nabla})$ denote the Deligne canonical extension of $(V\otimes \mathscr{O}_{\mathscr{C}^\circ}, \on{id}\otimes d)$ to $\mathscr{C}$. The Hodge filtration on $R^1\pi^\circ_*\mathbb{V}\otimes \mathscr{O}_{\mathscr{M}}$ is given, under the canonical identification $$R^1\pi^\circ_*\mathbb{V}\otimes \mathscr{O}_{\mathscr{M}}\overset{\sim}{\longrightarrow} R^1\pi_*(\mathscr{E}\overset{{\nabla}}{\longrightarrow}\mathscr{E}\otimes \Omega^1_{\mathscr{C}/\mathscr{M}}(\log \mathscr{D})),$$ by the filtration induced by the \emph{filtration b\^{e}te} on the relative de Rham complex $\mathscr{E}\overset{{\nabla}}{\longrightarrow}\mathscr{E}\otimes \Omega^1_{\mathscr{C}/\mathscr{M}}(\log \mathscr{D})$.
That is, $$F^1R^1\pi^\circ_*\mathbb{V}\otimes \mathscr{O}_\mathscr{M}:=\on{im}(\pi_*(\mathscr{E}\otimes \Omega^1_{\mathscr{C}/\mathscr{M}}(\log \mathscr{D}))\to R^1\pi^\circ_*\mathbb{V}\otimes \mathscr{O}_\mathscr{M}).$$
\end{theorem}
\begin{proof}
First, Griffiths transversality holds vacuously, as the Hodge filtration only has two steps.
To verify the remaining properties of a variation of mixed Hodge structure
we may do so pointwise.
Over a point,
the verification that
$R^1\pi^\circ_*\mathbb{V}_\mathbb{R}$ underlies a real variation of mixed Hodge
structure is carried out in
\cite[Theorem, p. 152]{timmerscheidt:mixed-hodge-structure-for-unitary}.
The description of the weight filtration follows from \cite[Lemma
6.2]{timmerscheidt:mixed-hodge-structure-for-unitary}.
The description of the Hodge filtration follows from the definition of the
Hodge filtration and the degeneration of the Hodge-de Rham spectral sequence
\cite[Theorem 7.1(a)]{timmerscheidt:mixed-hodge-structure-for-unitary}.

It remains only to verify graded-polarizability and admissibility. This follows from Saito's theory \cite{saito1990mixed} (also see
\cite{schnell2014overview}) as we now explain.
As $\mathbb{V}_\mathbb{R}$ is orthogonal, it underlies a real PVHS with trivial Hodge filtration and polarization arising from the orthogonal structure.
Hence, $Rj_*(\mathbb{V}_\mathbb{R})$ underlies an
object of the derived category of (real) mixed Hodge modules on $\mathscr{C}$. Hence
$R\pi^\circ_*\mathbb{V}_\mathbb{R}=R\pi_*Rj_*\mathbb{V}_\mathbb{R}$
underlies an object of the derived category of (real) mixed Hodge
modules on $\mathscr{M}$. As $\pi^\circ$ is a fiber bundle and hence $R^1\pi^\circ_*\mathbb{V}_\mathbb{R}$ is locally constant, this implies (by \cite[Theorem 21.1]{schnell2014overview}) that
$R^1\pi^\circ_*\mathbb{V}_\mathbb{R}$ in fact underlies a graded-polarizable
admissible variation of mixed Hodge structure. It remains only to compare it with
the structure described in the theorem. But this follows by compatibility
when $\mathscr{M}$ is a point, by base change.
\end{proof}
Let $\mathbb{V}$ be a unitary local system on $\mathscr{C}^\circ$, and let $\mathscr{H}_\mathbb{V}:=R^1\pi^\circ_*\mathbb{V}\otimes \mathscr{O}_\mathscr{M}$.   If $\mathbb{V}$ has a real structure $\mathbb{V}_\mathbb{R}$, as in \autoref{theorem:unitary-MHS}, then  $R^1\pi_*^\circ \mathbb{V}_\mathbb{R}$ underlies a natural polarized variation of $\mathbb{R}$-mixed Hodge structure. In general, we let $W^\bullet, F^\bullet$ be the weight and Hodge filtrations respectively, and $\overline{F}^\bullet$ the conjugate Hodge filtration induced by the natural isomorphism $$\mathscr{H}_{\mathbb{V}}\simeq \overline{\mathscr{H}_{\mathbb{V}^\vee}}.$$ Define $$\mathscr{H}_{\mathbb{V}}^{1,0}:= F^1\cap W^1,$$ $$\mathscr{H}^{0,1}_{\mathbb{V}} := \overline{F^1}\cap W^1,$$ and $$\mathscr{H}^{1,1}_{\mathbb{V}}=F^1\cap \overline{F^1}.$$
%If $\mathbb{V}$ is unitary but does not admit a real structure, we may still apply these constructions to $\mathbb{V}\oplus \mathbb{V}^{\vee}$, and define $$\mathscr{H}_\mathbb{V}^{i,j}:=\mathscr{H}_{\mathbb{V}\oplus \mathbb{V}^\vee}^{i,j}\cap \mathscr{H}_\mathbb{V}.$$
We give a schematic picture of the structures on $\mathscr{H}_{\mathbb{V}}$ in \autoref{figure:Hodge-diagram}.


\begin{remark}[Mixed Hodge structure]
\label{figure:Hodge-diagram}
\source{canonical-representations-surface-groups}
The source schematic organizes the weight and Hodge filtrations on
$\mathscr{H}_{\mathbb V}$ and the three summands introduced above.
\end{remark}
%

\begin{lemma}
\source{canonical-representations-surface-groups}
\notready
	\label{lemma:three-part-hodge-structure}
\source{canonical-representations-surface-groups}
	For $\mathbb V$ a unitary local system on $\mathscr C^\circ$, there is a natural
	direct sum decomposition
\begin{equation}
\label{eqn:hodge-decomposition-cpx}\mathscr{H}_\mathbb{V}=\mathscr{H}_{\mathbb{V}}^{1,0}\oplus
\source{canonical-representations-surface-groups}
\mathscr{H}_{\mathbb{V}}^{0,1}\oplus
\mathscr{H}_{\mathbb{V}}^{1,1}\end{equation}
Further, under the natural isomorphism
$\overline{\mathscr{H}_\mathbb{V}}=\mathscr{H}_{\mathbb{V}^\vee}$,
\begin{align}
	\label{equation:dual-conjugate}
\source{canonical-representations-surface-groups}
\overline{\mathscr{H}_{\mathbb{V}}^{i,j}}=\mathscr{H}_{\mathbb{V}^\vee}^{j,i}.
\end{align}
\end{lemma}
\begin{proof}
	One may check both statements hold
	by doing so on fibers, in which case it is a straightforward
	verification using the definitions of the weight and Hodge filtrations.
%	In this case, note that by construction
%$F^1=\mathscr{H}_{\mathbb{V}}^{1,0}\oplus \mathscr{H}_{\mathbb{V}}^{1,1}$ is a
%holomorphic subbundle of $\mathscr{H}_\mathbb{V}$.
\end{proof}

\subsection{Variations on the theorem of the fixed part}
\label{subsection:fixed-part}
\source{canonical-representations-surface-groups}

Our main goal here is to prove
\autoref{proposition:subrep-real-variation}, which is a variant of the theorem of the fixed part.
\begin{theorem}[Theorem of the fixed part, {\cite[4.19]{steenbrink1985variation}}, {\cite[14.52]{peters2008mixed}}]
\source{canonical-representations-surface-groups}
\notready\label{theorem:theorem-of-the-fixed-part}
	Let $\mathscr{X}$ be a smooth quasiprojective complex variety, and let $\mathbb{W}$ be an admissible graded-polarizable real variation of mixed Hodge structure on $\mathscr{X}$. Then there is a natural real Hodge structure on $H^0(\mathscr{X}, \mathbb{W})$ such that the natural inclusion $$H^0(\mathscr{X}, \mathbb{W})\otimes \underline{\mathbb{R}}\to \mathbb{W}$$ is a morphism of real variations of mixed Hodge structure.
\end{theorem}

%We next record a consequence to orient the reader, which we will not need in what follows.
%Using the direct sum decomposition as in
%\autoref{lemma:three-part-hodge-structure},
%we apply \autoref{theorem:theorem-of-the-fixed-part}
%to the MVHS associated to
%the constant sub-local system
%$H^0(\mathscr{M}, R^1\pi_*^\circ (\mathbb{V}\oplus \mathbb{V}^\vee)) \otimes \underline{\mathbb C} \subset
%R^1\pi_*^\circ (\mathbb{V}\oplus \mathbb{V}^\vee)$
%and obtain the following.
%\begin{corollary}\label{cor:unitary-thm-of-fixed-part}
%Let $\mathbb{V}$ be a unitary local system on $\mathscr{C}^\circ$, and let $\mathbb{H}_\mathbb{V}=H^0(\mathscr{M}, R^1\pi_*^\circ \mathbb{V})$.
%Then there is a natural direct sum decomposition $$\mathbb{H}_\mathbb{V}=\mathbb{H}_\mathbb{V}^{1,0}\oplus \mathbb{H}_\mathbb{V}^{0,1}\oplus \mathbb{H}_{\mathbb{V}}^{1,1}$$ such that the natural evaluation map $\mathbb{H}_\mathbb{V}\otimes \mathscr{O}_\mathscr{M}\to \mathscr{H}_\mathbb{V}$ sends $\mathbb{H}_\mathbb{V}^{i,j}\otimes \mathscr{O}_M$ into $\mathscr{H}_\mathbb{V}^{i,j}$ for all $i,j$.
%\end{corollary}

For our main results, will need a variant of the theorem of the fixed part for irreducible representations appearing
the cohomology of a unitary local system on a family of curves.
Given a complex local system $\mathbb L$ on a variety, we define the local system
$\widetilde{\mathbb L}$ by
\begin{align}
	\label{equation:real-construction}
\source{canonical-representations-surface-groups}
	\widetilde{\mathbb L} :=
	\begin{cases}
		\mathbb L & \text{ if $\mathbb L$ admits a real structure}  \\
		\mathbb L \oplus \overline{\mathbb L} & \text{ otherwise.}  \\
	\end{cases}
\end{align}
Note that $\widetilde{\mathbb L}$ admits a natural real structure. Note that if $\mathbb{L}$ is an irreducible local system underlying a complex variation of Hodge structure (necessarily unique up to reindexing), and if $\mathbb{L}$ admits a (necessarily unique) real structure, then $\mathbb{L}$ in fact underlies a real variation of Hodge structure. Indeed, the complex variation carried by $\mathbb{L}$ must be preserved (up to reindexing) by complex conjugation, as otherwise $\mathbb{L}$ would carry two distinct complex variations.
\begin{proposition}
\source{canonical-representations-surface-groups}
\notready\label{proposition:subrep-real-variation}
	Let $\mathscr{X}$ be a smooth quasiprojective complex variety, and let $\mathbb{W}$ be an admissible graded-polarizable real variation of mixed Hodge structure on $\mathscr{X}$. Let $\mathbb{L}$ be an irreducible complex local system on $\mathscr{X}$ such that $\Hom(\mathbb{L}, \mathbb{W}_\mathbb{C})\not = 0.$
	Then, $\widetilde{\mathbb L}$ underlies a polarizable real variation of Hodge structure, and there exists a nonzero \emph{constant} real  mixed Hodge structure $Q$ and a nonzero map of variations of mixed Hodge structure $$Q\otimes \widetilde{\mathbb{L}}\to \mathbb{W}.$$
	\end{proposition}
	\begin{proof}
	First, let $i$ be an integer for which there exists a nonzero map $\mathbb{L}\to
	\on{gr}^i_W\mathbb{W}_\mathbb{C}$. Then $$\on{gr}^i_W
	\mathbb{W}_\mathbb{C}=\oplus_j \mathbb{V}_j\otimes W_j,$$ where the
	$\mathbb{V}_j$ are distinct irreducible local systems carrying
	polarizable complex variations of Hodge structure, and the $W_j$ are
	constant complex variations \cite[Proposition 4.1.4(2)]{LL:geometric-local-systems}. By assumption $\mathbb{L}$
	occurs among the $\mathbb{V}_j$, and hence carries a complex PVHS.
	Hence
	$\widetilde{\mathbb{L}}$ carries a real PVHS.

		Let $Q:=\Hom(\widetilde{\mathbb{L}},
		\mathbb{W})=H^0(\mathscr{X}, \widetilde{\mathbb{L}}^\vee\otimes
		\mathbb{W})$. The real variation
		$\widetilde{\mathbb{L}}^\vee\otimes \mathbb{W}$ is a tensor
		product of a pure (hence admissible) variation with an
		admissible variation, hence is admissible itself. By the theorem
		of the fixed part, \autoref{theorem:theorem-of-the-fixed-part}, $Q$ carries a canonical real mixed Hodge structure such that the evaluation map $$Q\otimes \widetilde{\mathbb{L}}\to \mathbb{W}$$ is a map of variations, as desired. The map is nonzero by construction.
	\end{proof}

%	\begin{corollary}
%	\label{corollary:conjugation-trick}
%	Let $\mathbb{V}$ be a unitary local system on $\mathscr{C}^\circ$, and set $\mathbb{W}=R^1\pi_*^\circ\mathbb{V}$. Let $\mathbb{L}$ be an irreducible complex local system on $\mathscr{M}$ such that $\Hom(\mathbb{L}, \mathbb{W})\not=0$. Then $\mathbb{L}$ underlies a complex variation of Hodge structure, and there exists a bigraded vector space $Q$ and a nonzero flat map $$Q\otimes \mathbb{L}\otimes \mathscr{O}_{\mathscr{M}}\to \mathscr{H}_\mathbb{V}=\mathbb{W}\otimes \mathscr{O}_\mathscr{M}$$ respecting the grading on both sides.
%\end{corollary}
%\begin{proof}
%The statement is immediate from \autoref{proposition:subrep-real-variation},
%applied to $R^1\pi_*^\circ(\mathbb{V}\oplus \mathbb V^\vee)$.
%\end{proof}

\subsection{Consequences of non-abelian Hodge theory}
The main result from non-abelian Hodge theory that we will use is the following result of Mochizuki, generalizing earlier work of Simpson in the projective case:
\begin{theorem}[~\protect{\cite[Theorem 10.5]{mochizuki2006kobayashi}}]
\source{canonical-representations-surface-groups}
\notready
	\label{theorem:deformation-to-pvhs}
\source{canonical-representations-surface-groups}
	Let $\overline{X}$ be a smooth projective variety and $D\subset
	\overline{X}$ a strict normal crossings divisor. Let
	$X=\overline{X}\setminus D$.
	Any representation $$\rho: \pi_1(X)\to
	\on{GL}_r(\mathbb{C})$$ with finite determinant admits a
	deformation with constant determinant to a representation underlying a complex PVHS.
\end{theorem}
\begin{proof}
Aside from the statement about determinants, this is precisely	\cite[Theorem 10.5]{mochizuki2006kobayashi}; the statement about determinants follows by examining the proof of that theorem.
\end{proof}

We also use the related result that cohomologically rigid representations underlie complex PVHS.
\begin{lemma}
\source{canonical-representations-surface-groups}
\notready\label{lemma:rigid-reps}
Let $X, D, \rho$ be as in \autoref{theorem:deformation-to-pvhs}. Suppose in
addition that $\rho$ is semisimple with finite determinant and $H^1(X, \on{ad}\rho)=0$. Then $\rho$ underlies a complex PVHS.
\end{lemma}
\begin{proof}
	Let $G$ denote the Zariski-closure of the image of $\rho$.
	Note $G$ is reductive by semisimplicity of $\rho$. Let $\mathfrak{g}$ be
	the Lie algebra of $G$, viewed as a $\pi_1(X)$-representation. Note
	that $\mathfrak{g}$ is a subalgebra of $\on{ad}\rho$ as $\rho$ has
	finite determinant, and is moreover a summand of $\on{ad}\rho$ as
	$\on{ad}\rho$ is semisimple. Hence $\rho$ is cohomologically rigid as a
	$G$-representation, as $H^1(X, \mathfrak{g}) \subset H^1(X, \on{ad}
	\rho) =0$.

	We may apply \cite[Lemma 10.13]{mochizuki2006kobayashi},
	which yields that
	$\rho$ admits a deformation as a $G(\mathbb C)$-representation to
	$\rho_0: \pi_1(X) \to G(\mathbb C)$, which
	underlies a PVHS.
	But
	$\rho$ is rigid and hence admits no
	non-trivial deformations.	This implies $\rho = \rho_0$, so $\rho$ underlies a PVHS.
\end{proof}
